% GENERATED FILE. Edit canonical lessons, then run npm run generate:books.
% canonical-source-hash: fnv1a64:afd98bf2b748867d
% canonical-lessons: GU-C509-dhaman, GU-C509-phukvu, GU-C509-tipvu, GU-C509-kat, GU-C509-katavu

\chapter{Bellows, Hammer and Rust}
\label{ch:gu-bellows-hammer-and-rust}
\hlchaptermodality{509}

\begin{chapteropening}
\textbf{By the end of this chapter:} \emph{I can say how iron is heated, beaten and how it rusts.}

Complete the last lesson of chapter 509: I can say how iron is heated, beaten and how it rusts.
\end{chapteropening}

\section[\texorpdfstring{\gu{ધમણ} (dhamaṇ)}{ધમણ (dhamaṇ)}]{\gu{ધમણ} (dhamaṇ) — bellows}

\label{lesson:GU-C509-dhaman}



\begin{quote}
Before the new one: say the Gujarati for to heat (something), then the Gujarati for an anvil.
\end{quote}

\subsection*{You'll want to know: \gu{ધમણ}}
\textbf{\gu{ધમણ}} — \emph{dhamaṇ} — \textquotedblleft{}bellows\textquotedblright{}.

\begin{quote}
\gu{લુહારનો} \gu{દીકરો} \gu{ધમણ} \gu{ચલાવે} \gu{છે}.
\end{quote}

\begin{quote}
\emph{luhārno dīkro dhamaṇ calāve che.}
\end{quote}

\begin{quote}
The blacksmith's son works the bellows.
\end{quote}

A feminine noun.

\begin{cousinweb}
\textbf{\gu{ધમણ}} (\emph{dhamaṇ}) is inherited from Sanskrit \emph{dhamanī}, \textquotedblleft{}bellows\textquotedblright{}, from a root meaning \textquotedblleft{}to blow\textquotedblright{}, and CDIAL lists the Gujarati word (CDIAL 6734).
\end{cousinweb}

\begin{grammarlens}[title={What you've built}]
One more word for describing what you lived and felt.
\end{grammarlens}

\subsection*{Your turn}
\begin{itemize}
\raggedright
  \item \emph{Say it:} \emph{dhamaṇ}
  \item \emph{Say it:} \emph{dhamaṇ}, once more
  \item \emph{Say it:} \emph{luhārno dīkro dhamaṇ calāve che.} — The blacksmith's son works the bellows.
  \item \emph{Recall:} say the Gujarati for to heat (something), then the Gujarati for an anvil, then say \emph{dhamaṇ} again
\end{itemize}

\begin{tcolorbox}[breakable,colback=teal!4,colframe=teal!35!black,title={Before you move on}]
What does \gu{ધમણ} mean? (\textquotedblleft{}Bellows\textquotedblright{}.) Say it once more.
\end{tcolorbox}
\section[\texorpdfstring{\gu{ફૂંકવું} (phū̃kvũ)}{ફૂંકવું (phū̃kvũ)}]{\gu{ફૂંકવું} (phū̃kvũ) — to blow out (a flame)}

\label{lesson:GU-C509-phukvu}



\begin{quote}
Before the new one: say the Gujarati for an anvil, then the Gujarati for bellows.
\end{quote}

\subsection*{You'll want to know: \gu{ફૂંકવું}}
\textbf{\gu{ફૂંકવું}} — \emph{phū̃kvũ} — \textquotedblleft{}to blow out (a flame)\textquotedblright{}.

\begin{quote}
\gu{જન્મદિવસે} \gu{મીરાએ} \gu{બધી} \gu{મીણબત્તીઓ} \gu{ફૂંકી}.
\end{quote}

\begin{quote}
\emph{janmadivse Mīrāe badhī mīṇbattīo phū̃kī.}
\end{quote}

\begin{quote}
On her birthday Mira blew out all the candles.
\end{quote}

Past, with \gu{તેણે} (he or she …): \gu{તેણે} \gu{મીણબત્તી} \gu{ફૂંકી}, \emph{teṇe mīṇbattī phū̃kī}.

\begin{cousinweb}
\textbf{\gu{ફૂંકવું}} (\emph{phū̃kvũ}) comes from \emph{*phūtka}, \textquotedblleft{}blowing\textquotedblright{}, a form CDIAL reconstructs, and CDIAL lists the Gujarati verb (CDIAL 9102). The wind's own verb is \textbf{\gu{ફૂંકાવું}}.
\end{cousinweb}

\begin{grammarlens}[title={What you've built}]
One more word for describing what you lived and felt.
\end{grammarlens}

\subsection*{Your turn}
\begin{itemize}
\raggedright
  \item \emph{Say it:} \emph{phū̃kvũ}
  \item \emph{Say it:} \emph{phū̃kvũ}, once more
  \item \emph{Say it:} \emph{janmadivse Mīrāe badhī mīṇbattīo phū̃kī.} — On her birthday Mira blew out all the candles.
  \item \emph{Recall:} say the Gujarati for an anvil, then the Gujarati for bellows, then say \emph{phū̃kvũ} again
\end{itemize}

\begin{tcolorbox}[breakable,colback=teal!4,colframe=teal!35!black,title={Before you move on}]
What does \gu{ફૂંકવું} mean? (\textquotedblleft{}To blow out (a flame)\textquotedblright{}.) Say it once more.
\end{tcolorbox}
\section[\texorpdfstring{\gu{ટીપવું} (ṭīpvũ)}{ટીપવું (ṭīpvũ)}]{\gu{ટીપવું} (ṭīpvũ) — to beat (metal) flat}

\label{lesson:GU-C509-tipvu}



\begin{quote}
Before the new one: say the Gujarati for bellows, then the Gujarati for to blow out.
\end{quote}

\subsection*{You'll want to know: \gu{ટીપવું}}
\textbf{\gu{ટીપવું}} — \emph{ṭīpvũ} — \textquotedblleft{}to beat (metal) flat\textquotedblright{}.

\begin{quote}
\gu{લુહારે} \gu{ગરમ} \gu{લોખંડ} \gu{હથોડીથી} \gu{ટીપ્યું}.
\end{quote}

\begin{quote}
\emph{luhāre garam lokhanḍ hathoḍīthī ṭīpyũ.}
\end{quote}

\begin{quote}
The blacksmith beat the hot iron flat with the hammer.
\end{quote}

Past, with \gu{તેણે} (he or she …): \gu{તેણે} \gu{લોખંડ} \gu{ટીપ્યું}, \emph{teṇe lokhanḍ ṭīpyũ}.

\begin{cousinweb}
\textbf{\gu{ટીપવું}} (\emph{ṭīpvũ}) is said of metal, and of dough patted flat by hand. CDIAL lists a Gujarati \emph{ṭīpvũ} under a reconstructed \emph{*ṭipp}, \textquotedblleft{}squeeze, nip\textquotedblright{}, with another sense, and Hindi \emph{ṭīpnā}, \textquotedblleft{}to press\textquotedblright{}, beside it (CDIAL 5464); beating flat is not recorded there, so its history is not settled.
\end{cousinweb}

\begin{grammarlens}[title={What you've built}]
One more word for describing what you lived and felt.
\end{grammarlens}

\subsection*{Your turn}
\begin{itemize}
\raggedright
  \item \emph{Say it:} \emph{ṭīpvũ}
  \item \emph{Say it:} \emph{ṭīpvũ}, once more
  \item \emph{Say it:} \emph{luhāre garam lokhanḍ hathoḍīthī ṭīpyũ.} — The blacksmith beat the hot iron flat with the hammer.
  \item \emph{Recall:} say the Gujarati for bellows, then the Gujarati for to blow out, then say \emph{ṭīpvũ} again
\end{itemize}

\begin{tcolorbox}[breakable,colback=teal!4,colframe=teal!35!black,title={Before you move on}]
What does \gu{ટીપવું} mean? (\textquotedblleft{}To beat (metal) flat\textquotedblright{}.) Say it once more.
\end{tcolorbox}
\section[\texorpdfstring{\gu{કાટ} (kāṭ)}{કાટ (kāṭ)}]{\gu{કાટ} (kāṭ) — rust}

\label{lesson:GU-C509-kat}



\begin{quote}
Before the new one: say the Gujarati for to blow out, then the Gujarati for to beat (metal) flat.
\end{quote}

\subsection*{You'll want to know: \gu{કાટ}}
\textbf{\gu{કાટ}} — \emph{kāṭ} — \textquotedblleft{}rust\textquotedblright{}.

\begin{quote}
\gu{વરસાદ} \gu{પછી} \gu{લોખંડના} \gu{દરવાજાને} \gu{કાટ} \gu{લાગ્યો}.
\end{quote}

\begin{quote}
\emph{varsād pachī lokhanḍnā darvājāne kāṭ lāgyo.}
\end{quote}

\begin{quote}
After the rain the iron gate went rusty.
\end{quote}

A masculine noun.

\begin{cousinweb}
\textbf{\gu{કાટ}} (\emph{kāṭ}) is filed by CDIAL under Sanskrit \emph{kiṭṭa}, \textquotedblleft{}secretion, dross\textquotedblright{}, and CDIAL lists the Gujarati word as \textquotedblleft{}rust\textquotedblright{} (CDIAL 3156). Rust strikes the metal: \textbf{\gu{દરવાજાને} \gu{કાટ} \gu{લાગ્યો}}.
\end{cousinweb}

\begin{grammarlens}[title={What you've built}]
One more word for describing what you lived and felt.
\end{grammarlens}

\subsection*{Your turn}
\begin{itemize}
\raggedright
  \item \emph{Say it:} \emph{kāṭ}
  \item \emph{Say it:} \emph{kāṭ}, once more
  \item \emph{Say it:} \emph{varsād pachī lokhanḍnā darvājāne kāṭ lāgyo.} — After the rain the iron gate went rusty.
  \item \emph{Recall:} say the Gujarati for to blow out, then the Gujarati for to beat (metal) flat, then say \emph{kāṭ} again
\end{itemize}

\begin{tcolorbox}[breakable,colback=teal!4,colframe=teal!35!black,title={Before you move on}]
What does \gu{કાટ} mean? (\textquotedblleft{}Rust\textquotedblright{}.) Say it once more.
\end{tcolorbox}
\section[\texorpdfstring{\gu{કટાવું} (kaṭāvũ)}{કટાવું (kaṭāvũ)}]{\gu{કટાવું} (kaṭāvũ) — to rust}

\label{lesson:GU-C509-katavu}



\begin{quote}
Before the new one: say the Gujarati for to beat (metal) flat, then the Gujarati for rust.
\end{quote}

\subsection*{You'll want to know: \gu{કટાવું}}
\textbf{\gu{કટાવું}} — \emph{kaṭāvũ} — \textquotedblleft{}to rust\textquotedblright{}.

\begin{quote}
\gu{વરસાદમાં} \gu{બધાં} \gu{ઓજારો} \gu{કટાઈ} \gu{ગયાં}.
\end{quote}

\begin{quote}
\emph{varsādmā̃ badhā̃ ojāro kaṭāī gayā̃.}
\end{quote}

\begin{quote}
All the tools rusted in the rain.
\end{quote}

Past (it …): \gu{તે} \gu{છરી} \gu{કટાઈ} \gu{ગઈ}, \emph{te charī kaṭāī gaī}.

\begin{cousinweb}
\textbf{\gu{કટાવું}} (\emph{kaṭāvũ}) is built on \textbf{\gu{કાટ}}, \textquotedblleft{}rust\textquotedblright{}, which CDIAL files under Sanskrit \emph{kiṭṭa}, \textquotedblleft{}dross\textquotedblright{} (CDIAL 3156). It is intransitive: \textbf{\gu{છરી} \gu{કટાઈ} \gu{ગઈ}}.
\end{cousinweb}

\begin{grammarlens}[title={What you've built}]
One more word for describing what you lived and felt.
\end{grammarlens}

\subsection*{Your turn}
\begin{itemize}
\raggedright
  \item \emph{Say it:} \emph{kaṭāvũ}
  \item \emph{Say it:} \emph{kaṭāvũ}, once more
  \item \emph{Say it:} \emph{varsādmā̃ badhā̃ ojāro kaṭāī gayā̃.} — All the tools rusted in the rain.
  \item \emph{Recall:} say the Gujarati for to beat (metal) flat, then the Gujarati for rust, then say \emph{kaṭāvũ} again
\end{itemize}

\begin{tcolorbox}[breakable,colback=teal!4,colframe=teal!35!black,title={Before you move on}]
What does \gu{કટાવું} mean? (\textquotedblleft{}To rust\textquotedblright{}.) Say it once more.
\end{tcolorbox}
